\documentclass[12pt]{beamer}

% ****************
% ***** INFO *****
% ****************
\usepackage[T2A]{fontenc}
\usepackage[russian]{babel}
\usepackage[utf8]{inputenc}
\usepackage{amsmath, amssymb}
\usepackage{geometry, csquotes}
\usepackage[style = ieee]{biblatex}
\addbibresource{refs_1.bib}
\usetheme{Madrid}

\def\epsimage#1#2#3#4{
  \begin{figure}[!ht]
    \begin{center}%
      \includegraphics[height=#2\textheight]{#1.eps}\caption{#4}\label{#3}%
    \end{center}
\end{figure}}
\renewcommand*{\bibfont}{\normalfont\footnotesize}
% \usepackage{helvet}
\title[Расстояние Г-Х между облаками...]{Расстояние Громова--Хаусдорфа
  между облаками с нетривиальным
пересечением стационарных групп}
\author{Нестеров Б. А.}
\institute{607 группа}
\date{2026}

\begin{document}

\frame{\titlepage}

\begin{frame}{Основные определения}
  \begin{itemize}
    \item \textbf{Симплекс} --- метрическое пространство \( \Delta _ \alpha
      \), в котором все расстояния равны \( 1 \).
    \item \textbf{Расстояние Громова–Хаусдорфа} $d_{GH}$ --- метрика на
      классах изометричных компактных пространств. Для неограниченных
      пространств $d_{GH}$ может быть \( 0 \) и $\infty$.
    \item \textbf{Соответствие} между пространствами \( X, Y \) ---
      подмножество \( X \times Y \), проекции которого на \( X \) и
      \( Y \) сюръективны.
  \end{itemize}
\end{frame}
\begin{frame}{Основные определения}
  \begin{itemize}
    \item \textbf{Облако} $[X]$ --- класс эквивалентности: $Y\sim X \iff
      d_{GH}(X,Y)<\infty$.
    \item \textbf{Стационарная группа} $\mathrm{St}\big([X]\big) =
      \bigl\{\lambda>0
      : [\lambda X]=[X]\bigr\}$.
    \item \textbf{Центр} облака с нетривиальной стац. группой --- единственное
      $M$: $\lambda M = M$ для $\lambda\in\mathrm{St}\big([M]\big)$.
  \end{itemize}
\end{frame}

\begin{frame}{Облака --- собственные классы}
  \begin{block}{Лемма}
    В любом облаке лежат пространства сколь угодно большой мощности.
  \end{block}
  \pause
  \begin{block}{Теорема 3.3}
    Все облака являются собственными классами в системе аксиом NBG.
  \end{block}
\end{frame}

\begin{frame}{Расстояние между облаками с пересечением стац. групп}
  \begin{block}{Лемма 3.6}
    Пусть $\mathrm{St}([X])\cap\mathrm{St}([Y])\neq\{1\}$. Тогда
    $d_{GH}([X],[Y])$ равно $0$ или $\infty$.
  \end{block}
  \pause
  \begin{block}{Примеры}
    Облака \( [\mathbb{R}^n], n \in \mathbb{N} \), \( [\Delta _1] \)
    все попарно имеют
    нетривиальное пересечение стационарных групп, равное \( \mathbb{R}^+ \).
  \end{block}
\end{frame}

\begin{frame}{Теорема об образе центра}
  \begin{block}{Теорема 4.3}
    Пусть $[M]$ --- облако с нетривиальной стац. группой, $M$ --- его
    центр. $R$ --- соответствие между $[\Delta_1]$ и $[M]$ с искажением
    $\varepsilon$. Тогда образ $\Delta_1$ лежит на расстоянии $\le
    2\varepsilon$ от $M$.
  \end{block}
  \pause
  \textbf{Следствие (теорема 5.1):}
  Если в $[M]$ есть $Y_1,Y_2$ с $\max|Y_i,M|=r$, $|Y_1,Y_2|>r$, то
  $d_{GH}([\Delta_1],[M])=\infty$.
\end{frame}
\begin{frame}
  \epsimage{diffAngles}{0.50}{fig:diffAngles}{Иллюстрация
  несовпадения расстояний}
\end{frame}

\begin{frame}{Облако вещественной прямой}
  \begin{itemize}
    \item $d_{GH}(\mathbb{Z},\mathbb{R}) = \frac12$,
      $d_{GH}(\tilde{\mathbb{R}},\mathbb{R})\le\frac12$,
      $d_{GH}(\mathbb{Z},\tilde{\mathbb{R}})\ge\frac23$ (теорема 6.3).
    \item Пространства $\mathbb{Z}$ и $\tilde{\mathbb{R}}$
      удовлетворяют условиям теоремы 5.1.
    \item \textbf{Следствие 6.4:} $d_{GH}([\Delta_1],[\mathbb{R}]) = \infty$.
  \end{itemize}
  \pause
  \vfill
  \textbf{Замечание:} В облаке $[\Delta_1]$ выполняется
  ультраметрическое неравенство, в $[\mathbb{R}]$ --- нет.
\end{frame}

\begin{frame}{Углы облаков}
  \begin{block}{Определение}
    Угол между $X_1,X_2\in[M]$ ($M$ --- центр):
    \[
      \varphi(X_1,X_2)=\arccos\frac{r_1^2+r_2^2-d^2}{2r_1r_2},\quad
      r_i=|X_iM|,\ d=|X_1X_2|.
    \]
  \end{block}
  \pause
  \begin{block}{Определение}
    Угол облака $[M]$: $\varphi([M])=\sup\varphi(X_1,X_2).$

    Равнобедренный угол $\varphi_e([M])$ --- супремум по $|X_1M|=|X_2M|$.
  \end{block}
  \pause
  \textbf{Примеры:}
  \[
    \varphi([\Delta_1])=\frac\pi2,\quad
    \varphi_e([\Delta_1])=\frac\pi3,\quad
    \varphi([\mathbb{R}])=\varphi_e([\mathbb{R}])=\pi.
  \]
\end{frame}

\begin{frame}{Теорема об облаках с различными углами}
  \begin{block}{Теорема 7.9}
    Пусть $\mathrm{St}([M])\cap\mathrm{St}([N])\neq\{1\}$ и
    $\varphi([M])\neq\varphi([N])$. Тогда $d_{GH}([M],[N])=\infty$.
  \end{block}
\end{frame}

\begin{frame}{Выводы}
  \begin{itemize}
    \item Облака --- собственные классы, а не множества (в NBG).
    \item Расстояние между облаками с нетривиальным пересечением
      стационарных групп может быть только $0$ или $\infty$.
    \item Доказана теорема об образе центра и получено условие
      бесконечности расстояния до $[\Delta_1]$.
    \item Исследовано облако вещественной прямой: вычислены расстояния
      до подпространств.
    \item Введены новые инварианты --- угол и равнобедренный угол облака.
    \item Различие углов при нетривиальном пересечении стац. групп
      влечёт бесконечность расстояния Громова–Хаусдорфа.
  \end{itemize}
\end{frame}
\nocite{*}
\printbibliography
\end{document}
